\documentclass[10pt,a4paper]{amsart}
\usepackage[margin=19mm]{geometry}
\usepackage{amsmath,amssymb,mathtools}
\usepackage[scr=rsfs,cal=euler]{mathalfa}
\usepackage{xfrac}
\usepackage[unicode,hidelinks,bookmarks=false]{hyperref}
\newcommand{\ie}{i.e.\ }
\newcommand{\cf}{cf.\ }
\newcommand{\resp}{respectively\ }
\newcommand{\Subsection}{\subsection}
\providecommand{\nobreakdash}{\nobreak\hbox{-}}
\setlength{\emergencystretch}{2em}
\multlinegap0pt
\usepackage{mathtools}
\newcommand{\sq}[1]{\overline{#1^2}}
\theoremstyle{plain}
\newtheorem{theorem}{Theorem}[section]
\newtheorem{proposition}[theorem]{Proposition}
\newtheorem{lemma}[theorem]{Lemma}
\newtheorem{definition}[theorem]{Definition}
\newtheorem{corollary}[theorem]{Corollary}
\newtheorem{problem}[theorem]{Problem}
\newtheorem{algorithm}[theorem]{Algorithm}
\newtheorem{observation}[theorem]{Observation}
\newtheorem{example}{Example}[section]
\newtheorem{claim}{Claim}
\newtheorem{question}{Question}
\newtheorem{conjecture}{Conjecture}
\newtheorem{remark}{Remark}
\title[There are no nontrivial chordal square-complementary graphs]{There are no nontrivial chordal square-complementary graphs}
\author{Pierre Aboulker, Martin Milani\v{c}, Peter Mur\v{s}i\v{c}, Miguel A. Piza\~na}
\date{Source: arXiv:2609.28702v1 (CC BY 4.0)}
\begin{document}
\maketitle

\noindent\emph{Source and licence.} The delivered work is the article shown in the title above, version arXiv:2609.28702v1 (CC BY 4.0), distributed under the Creative Commons Attribution 4.0 International licence, as recorded in the arXiv licence metadata for this exact version and shown on the abstract page saved with this item. The full licence notice, which carries the licence URL and the address of that abstract page, is the bundled file \texttt{licenses/arxiv-2609.28702-CC-BY-4.0.txt}.
\par\smallskip

\noindent\textit{Editorial scope.} Counted unit: the article's Theorem (source label thm:notwins) (source0.tex lines 219--223, source label \texttt{thm:notwins}): ``Let be a minimal nontrivial squco graph. Then contains no (true or false) twin vertices. In particular, there is no nontrivial distance-hereditary squco graph.''. It is delivered together with the article's complete original proof (source0.tex lines 225--268, one proof environment adjacent to the statement, 44 lines), copied line for line from the article's own text. Required context retained uncounted: thm:notruetwins (lines 205--207). Editorial formatting only: an AMS \texttt{amsart} preamble that restates the article's own macro definitions and its own theorem declarations, and the theorem environments the retained lines use; the packages mathtools are loaded to match the article's own setup; the article's own class file is neither shipped nor input; the retained environments are renumbered sequentially in order of appearance, whereas the article prints them with its own numbering; the article's equation numbering is not reproduced and every cross-reference inside the retained text resolves to the wrapper's numbering. Exactly 1 bibliography entry is delivered, transcribed from the article's own thebibliography; All retained bodies are exact copies of the bundled \texttt{sources/211612/source0.tex}, so no omission receipt applies. No asset-inclusion command occurs inside any retained span and the package ships no image asset.


\section*{Retained context: thm:notruetwins (source0.tex lines 205--207)}

\begin{theorem}\label{thm:notruetwins}
  If $G$ is nontrivial and squco then $G$ does not contain dominated vertices, true twins, nor simplicial vertices. \qed
\end{theorem}


\section{The counted unit: Theorem (source label thm:notwins) and its complete original proof}

\begin{theorem}\label{thm:notwins}
  Let $G$ be a minimal nontrivial squco graph. 
  Then $G$ contains no (true or false) twin vertices. 
  In particular, there is no nontrivial distance-hereditary squco graph.
\end{theorem}

\begin{proof}
  Let $\varphi\colon \sq{G}\rightarrow G$ be an isomorphism viewed as a permutation of $V(G)$. 
  By Theorem~\ref{thm:notruetwins}, any pair of twins in $G$ must be false twins.  
  Suppose $G$ contains false twins.  
  Let $x\in V(G)$ be a twin, and let $T_1$ be the set of all false twins of $x$ (including $x$).  
  Let $\{T_{1}, T_{2}, \ldots, T_{s}\}$ be the $\varphi$-orbit of $T_{1}$. 
  We may assume that $T_{i} = \varphi^{i-1}(T_{1})$ for
  $i=1,2,\ldots, s$, and that $\varphi^{s}(T_{1}) = T_{1}$, with $s$ being the minimum such positive integer.  
  Since $\varphi$ is an isomorphism, each of these sets $T_{i}$ is a maximal set of pairwise twin vertices.
  Furthermore, since the relation of being false twins is an
  equivalence relation, these sets are pairwise disjoint. 
  Moreover, $|T_{i}|\geq 2$ for all $i$.

  Let $y=\varphi^s(x)$. 
  Suppose first that $y\neq x$.  
  Since $\varphi^{s}(T_{1}) = T_{1}$, we have $y\in T_{1}$, so $y$ is a twin of $x$.  
  Hence, the transposition $\alpha = (x\,\,\, y)$ is an automorphism of $G$.  
  It follows that $\psi = \alpha\circ \varphi \colon\sq{G} \rightarrow G$ is also an isomorphism. 
  Since $\alpha$ only exchanges $x$ and $y$, we have that $\psi(z)$ is identical to $\varphi(z)$ except when $\varphi(z)\in \{x,y\}$. 
  It follows that $\psi^{i}(x) = \varphi^{i}(x)$ for $i=0,1,\ldots, s-1$ and that $\psi^{s}(x) = \psi(\psi^{s-1}(x)) = \psi(\varphi^{s-1}(x)) = \alpha(\varphi(\varphi^{{s-1}}(x))) = \alpha(\varphi^{s}(x)) = \alpha(y) = x$. 
  Hence, the $\psi$-orbit of $x$ is exactly $\{x, \psi(x), \psi^2(x), \ldots, \psi^{s-1}(x)\}$.  
  Therefore, we may assume without loss of generality that $y=x$ and that the $\varphi$-orbit of $x$ is exactly
  $X \coloneqq\{x, \varphi(x), \varphi^2(x), \ldots, \varphi^{s-1}(x)\}$.

  Now, since each vertex in $X$ has a twin outside $X$, it follows that $G-X$ is connected and that the distances in $G-X$ are the same as in $G$ for every pair of vertices $x,y\in V(G) \setminus X$. 
  Let us show that the restriction of $\varphi$ to $V(G)\setminus X$ is an isomorphism from $\sq{(G-X)}$ to $G-X$. 
  Note first that $V\left(\sq{(G-X)}\right)=V(G) \setminus X =V(G-X)$. Since $X$ is $\varphi$-invariant, so is $V(G)-X$ and it follows that  $\varphi\left(V\left(\sq{(G-X)}\right)\right)=V(G-X)$. 
  Now take $x,y\in V(G) \setminus X$. 
  Then $xy\in E\left(\sq{(G-X)}\right)$ $\iff$
    $d_{G-X}(x,y)\geq 3$ $\iff$ $d_{G}(x,y)\geq 3$ $\iff$
        $xy\in E(\sq{G})$ $\iff$ $\varphi(x)\varphi(y)\in E(G)$ $\iff$
            $\varphi(x)\varphi(y)\in E(G-X)$.  
        Hence, $\sq{(G-X)}\cong G-X$, that is, $G-X$ is squco.  
        By the minimality of $G$, the graph $G-X$ must be trivial.  
        Recall that $T_{1}, T_{2},\ldots, T_{s}$ are pairwise disjoint, and that for all $i$, $|T_{i}|\geq 2$, but $|T_{i}\cap X| = 1$.  
            It follows that $s=1$ and hence, $G$            must have exactly two vertices, but there is no squco graph on two vertices.

  Finally, suppose that there exists a nontrivial distance-hereditary squco graph~$G_0$. Then, $G_0$ contains, as an induced subgraph, some minimal nontrivial squco graph $G$. 
Since the class of distance-hereditary graphs is closed under vertex deletions, $G$ is distance-hereditary. 
  Being distance-hereditary, it is well known that $G$
  contains either pendant vertices (which is incompatible with squco graphs; see~\cite{BST94}) or twins. 
  If $G$ contains twins, removing an orbit of twins as before would produce a smaller distance-hereditary graph $G-X$ which is also a squco graph. 
  As before, it follows that $G-X$ is trivial and that $G$ has exactly two vertices, a contradiction.
\end{proof}

\begin{thebibliography}{99}
\bibitem[BST94]{BST94} V. Balti\'{c}, S. Simi\'{c} and V. Tintor. \newblock {\em Some remarks on graph equation {$G^2=\overline G$}}. \newblock Univ. Beograd. Publ. Elektrotehn. Fak. Ser. Mat. {\bf 5} (1994) 43--48.
\end{thebibliography}
\end{document}
